\documentclass{proposal}

\usepackage[utf8]{inputenc}
\usepackage{hyperref}
\usepackage{wasysym}
\usepackage{amssymb}
\usepackage{tcolorbox}

\type{Proposal for a Semester project}
\title{\!\!\!\!A Language for Formalizing Legal Provisions\!\!\!\!}
\proposalsupervisor[Supervisors]{François Hublet, Jakob Merane (D-GESS)}
\proposalissuedate{February 2024}

\begin{document}

\vspace*{-3ex}
\section{Prerequisites}

\begin{compactitem}
\item Familiarity with runtime monitoring from courses (InfSec, FMSec) or self-study. 
\item Familiarity with compilers (e.g., from the Compiler design course).
\item Good OCaml programming skills.
\end{compactitem}

\section{Introduction}

While Privacy by Design (PbD)~\cite{Cavoukian2009} is prescribed by modern privacy
regulations such as the EU's GDPR, achieving PbD in real
software systems is a notoriously difficult task. One emerging
technique to realize PbD is \emph{Runtime enforcement} (RE)~\cite{Schneider2000}. In RE, an \emph{enforcer}, loaded with a specification of a system's privacy requirements, observes the actions performed by the system and instructs it to perform actions that will ensure compliance with these requirements at all times.
In our previous work~\cite{Hublet2023,Hublet2024}, we have shown how RE can be used to enforce parts of privacy regulations in real-world software systems. As a specification language, we used Metric First-Order Temporal Logic (MFOTL)~\cite{Chomicki1995,Basin2015}, an expressive temporal logic for which state-of-the-art enforcement tools exist~\cite{Hublet2022,Hublet2024}.
However, using MFOTL as a specification language presents two main limitations. First, MFOTL is not easily understandable by legal experts. Second, the MFOTL syntax does not provide native support for typical features of legal texts (e.g., default cases and exceptions).
To allow for the development of larger formal specifications of laws that can be validated with legal experts, a more user-friendly specification language is needed. If this language can be compiled into MFOTL, the enforcement tools developed in previous work~\cite{Hublet2022,Hublet2024} can be re-used as backends.

While an active field of (mostly theoretical) research exists
at the interplay of logic, linguistics, computer science, and law~\cite{Gabbay2010}, frameworks supporting the development of specifications of legal provisions amenable
to formal reasoning are still very rare. Recently, Merigoux et al. proposed \textsc{Catala}, a special-purpose programming language tailored for the formalization of tax codes~\cite{Merigoux2021,Huttner2022}. The complexity of tax codes lies in
complex conditional structures depending on a large number of variables and 
conditions, requiring the use of special (typically, default~\cite{Antoniou1999}) logics. 
However, tax codes lack most of the temporal and system-specific
dimensions that are intrinsic to privacy regulations.
Hence, Merigoux et al.'s language cannot be readily extended
to develop enforceable specifications of (especially privacy) laws.


\section{Objectives}\label{sec:objectives}

In this project, we will develop a new specification language and some tools that will allow technical and legal experts to create enforceable specifications of legal provisions. The final artifact should consist of:
\begin{enumerate}
\item A concise but complete user documentation of the new language's  features.
\item A compiler from our specification language into MFOTL, including a type-checker that `lifts' the type system proposed by Hublet, Lima, et al.~\cite{Hublet2024} to our new language (see details below).
\item A minimal set of tools for \href{https://code.visualstudio.com/}{VSCode} integration.
\end{enumerate}
The project will be co-supervised by members of the Institute of Information Security (François Hublet) and the Center for Law and Economics (Jakob Merane) at ETH Zurich. This joint supervision aims at ensuring that the specification language fits the needs of both legal and technical experts in terms of user-friendliness and choice of features.

In the rest of this section, we review the main expected features of our specification language.
\begin{description}
\item[Module system.] The language supports the definition of modules, import of modules from other files and folders, and inheritance between modules.
\item[Structure of legal texts.] {\color{green} The language supports a system of labels to decompose the specification into articles, paragraphs, points...}
\item[Description of ontology / signature.] {\color{green} The language supports the definition of system ontologies that correspond to \emph{signatures} in RE. To this end, users write a series of \emph{event declarations}. Every event has a name and a sequence of argument names and types. Additionally, users can specify a high-level English semantics for each event.}

\emph{Optional}: Add support for record / algebraic types in arguments of events~\cite{LimaGraf2023}. Add support for events that capture a state persisted between time points, e.g., from a permissions database. Add variable-style syntax, e.g., $\mathtt{limit=2}$ instead of $\mathtt{Limit}(2)$, for unary events.
\item[Types of events.] {\color{green} Each event can be marked as `observable' (= the enforcer can observe the corresponding system action), `suppressable' (= the enforcer can suppress the corresponding system action), `causable' (= the enforcer can cause the corresponding system action), or a combination of these. Additionally, events can be marked as `internal' to denote that they can only be generated by the enforcer for internal bookkeeping.}
\item[Rules.] {\color{green} Users can specify rules of the form $\phi \Longrightarrow \psi$ than can encode obligations, permissions, and constitutive rules (see~\cite{Robaldo2020} for this distinction). Each rule can be accompanied by a comment in English.} The language must support full MFOTL syntax, but also provide basic temporal patterns out of the box. E.g., instead of $\Box (\phi \Longrightarrow \lozenge_{[0,10d]} \psi)$, users can write \texttt{whenever} $\phi$ \texttt{then within 10 days ensure} $\psi$.
\item[Type system.] {\color{green} Users can mark rules as `enforceable'.} An algorithm analogous to the one presented by Hublet, Lima, et al.~\cite{Hublet2024} checks that these markings are consistent with the types of the events and returns errors if this is not the case.
\item[Exceptions.] {\color{green} Rules can be specified to introduce exceptions to previously stated rules, in the style of default logic.} This can be easily compiled into MFOTL by generate special exception events that are reinjected into the LHS of the overwritten rules.

\emph{Optional:} Add support for rule templates to be able to directly encode provisions such as, e.g., `if [some condition] then rule [some rule] is applicable, but the delay of ten days in [some rule] is replaced by a delay of thirty days'.
\item[Explicit reasoning about obligations / permissions.] Users can express rules of the form `if we ought (are allowed) to do $\phi$, then we ought (are allowed) to do $\psi$'.

\item[Support for refinement.] Users can declare \emph{refinements} of ontologies in which every event from an abstract ontology is mapped to a formula from the concrete ontology. The type-checker checks that the refinement preserves the types of the events.

\item[Compilation.] A specification is compiled into a signature with some causability/suppressability marks~\cite{Hublet2024} and a formula or set of formulae to be enforced by WhyEnf~\cite{Hublet2024} or monitored by WhyMon~\cite{Lima2023}.  The user can chose whether to compile all rules or only the enforceable ones. 

\emph{Optional:} Allow users to automatically generate documentation for their specification in the form of a set of HTML pages.
\end{description}

\section{Tasks}\label{sec:tasks}

The project is decomposed into the following tasks:

\begin{enumerate}
\item Familiarize yourself with the WhyEnf tool for runtime enforcement of MFOTL policies and its type system~\cite{Hublet2024}. Access to the source code will be provided to you by the start of the project. Study the most closely related works~\cite{Robaldo2020,Merigoux2021,Hublet2023}.
\item Implement a parser, type-checker, and compiler for the new specification language in OCaml. We will proceed feature by feature, designing the syntax of the language in an iterative manner. Provide a concise documentation of each language feature, with some examples. As you progress, also create some basic syntax highlighting for VSCode.
\item Create a VSCode \href{https://code.visualstudio.com/api/language-extensions/overview}{extension} for our language to provide seemless iteraction with the type-checker and compiler.
\item Prepare a demonstration on previously formalized GDPR articles (see, e.g., \cite{Kvamme2024}).
\end{enumerate}

\section{Deliverables}\label{sec:deliverables}


The following deliverables are due at the end of the project:
\begin{itemize}
    \item Full development (compiler, VSCode extension).
    \item Technical documentation of the language, typeset in English using LaTeX (10-30 pages with examples).
    \item A 20-minute talk giving an overview of the key findings of the project, including a short demo.
\end{itemize}


{%\setlength\bibsep{2pt}
\small
\bibliographystyle{unsrt}
\bibliography{flows}
}

\end{document}

%%% Local Variables:
%%% mode: latex
%%% TeX-master: t
%%% End:
